\begin{proof}[Proof of \cref{obj:thm:supported-boundaries}]
\begin{enumerate}
\item Choose the universal constants
\[
c:=\frac{1}{2560000\pi^2},
\qquad
C_0:=21.
\]
Both are positive. Fix admissible \(n,d,q\) and a joint assignment--audit law satisfying the stated assumptions. Independence and the prescribed marginals give
\[
\operatorname{Law}(Z,W)
=
\left(\bigotimes_{j\in V}\operatorname{Bernoulli}(1/2)\right)
\otimes
\left(\bigotimes_{(j,i)\in V^2:j\ne i}
\operatorname{Bernoulli}(q)\right).
\]
Thus the experiment has the independent assignment--audit product design. % lean: CausalSmith.Experimentation.ThinnedgraphAdditiveRiskFrontier.design_eq_thinnedDesign
Applying \cref{obj:thm:observable-upper}, with the envelope defined in \cref{obj:def:upper-envelope}, yields
\[
R_{n,d}(q)\le \min\{1,u(n,d,q)\}.
\]
% lean: CausalSmith.Experimentation.ThinnedgraphAdditiveRiskFrontier.observable_upper

\item By \cref{obj:lem:supplied-block-record-lower},
\[
R_{n,d}(q)\ge
\frac{1}{160000\pi^2}
\min\left\{1,\frac{(d+1)^2}{n}\right\}.
\]
Since
\[
0<c=\frac{1}{2560000\pi^2}
\le \frac{1}{160000\pi^2},
\qquad
\min\left\{1,\frac{(d+1)^2}{n}\right\}\ge0,
\]
we obtain
\[
R_{n,d}(q)\ge
c\min\left\{1,\frac{(d+1)^2}{n}\right\}.
\]
% lean: CausalSmith.Experimentation.ThinnedgraphAdditiveRiskFrontier.supported_boundaries

\item We next establish the degree-one frontier bound
\[
R_{n,1}(q)\ge cF(n,1,q),
\]
where the risk scale of \cref{obj:def:frontier-scale} specializes to
\[
F(n,1,q)=
\begin{cases}
\min\{1,1/(nq)\},&q>0,\\
1,&q=0.
\end{cases}
\qquad
0<F(n,1,q)\le1.
\]
The auxiliary degree \(1\) is admissible because \(n\ge4\).

First suppose \(1/n\ge1/16\). Then
\[
\frac4n\ge\frac1{16},
\qquad
\min\left\{1,\frac4n\right\}\ge\frac1{16}.
\]
Using \cref{obj:lem:supplied-block-record-lower} at degree one and
\(2560000=16\cdot160000\), we conclude
\[
cF(n,1,q)
\le c
\le
\frac{1}{160000\pi^2}\min\left\{1,\frac4n\right\}
\le R_{n,1}(q).
\]

Now suppose \(1/n<1/16\). Then \(n>16\), and in particular \(4<n\). Define the integer block count and source count by
\[
B:=\left\lfloor\frac n2\right\rfloor,
\qquad
m:=B.
\]
The division remainder satisfies
\[
0\le n-2B<2,
\qquad
2m\le n<2m+2.
\]
Consequently \(B\ge1\), and, since \(2<n/2\),
\[
n<2m+2<2m+\frac n2
\quad\Longrightarrow\quad
\frac n4\le m.
\]
Thus
\[
B\ge1,\qquad 2m\le n,\qquad \frac n4\le m\le n.
\]
% lean: CausalSmith.Experimentation.ThinnedgraphAdditiveRiskFrontier.frontier_block_rounding

Define the degree-one testing scale and testing constant by
\[
s:=s(B,1,q)=
\begin{cases}
\min\{1,1/\sqrt{mq}\},&q>0,\\
1,&q=0,
\end{cases}
\qquad
\kappa:=\frac1{100\pi}.
\]
These are the quantities in \cref{obj:thm:block-testing-scale}, since the association-revelation probability at degree one is
\[
p:=1-(1-q)^1=q.
\]
For \(q>0\), both entries in the minimum defining \(s\) are positive; for \(q=0\), \(s=1\). Hence
\[
0<s\le1,\qquad \kappa>0.
\]
% lean: CausalSmith.Experimentation.ThinnedgraphAdditiveRiskFrontier.testingScale_pos_le_one

We claim that
\[
F(n,1,q)\le s^2.
\]
For \(q>0\), the positive denominators and \(m\le n\) imply
\[
\left(\frac1{\sqrt{mq}}\right)^2=\frac1{mq},
\qquad
\frac1{nq}\le\frac1{mq}.
\]
If \(1/\sqrt{mq}\le1\), then
\[
F(n,1,q)\le\frac1{nq}
\le\frac1{mq}=s^2.
\]
If \(1/\sqrt{mq}>1\), then
\[
F(n,1,q)\le1=s^2.
\]
At \(q=0\), the definitions give
\[
F(n,1,0)=1=s^2.
\]
This proves the claim over the entire retention range.
% lean: CausalSmith.Experimentation.ThinnedgraphAdditiveRiskFrontier.frontierScale_le_testingScale_sq

Set the testing amplitude and target magnitude to
\[
h:=\kappa s,
\qquad
\Delta:=\frac{mh}{n}.
\]
Since \(\pi>3\),
\[
0<h\le\kappa=\frac1{100\pi}<\frac14,
\qquad
\Delta>0.
\]
Take the opposite-sign block priors of \cref{obj:def:block-prior} with \(B\) blocks, degree one, and amplitude \(h\). The conditions \(B\ge1\), \(2B\le n\), and \(0<h\le1/4\) allow us to apply \cref{obj:lem:block-law}. It gives, under the respective priors,
\[
\theta\in\mathcal M_{n,1},
\qquad
\tau(\theta)=\Delta
\quad\text{and}\quad
\tau(\theta)=-\Delta
\quad\text{almost surely}.
\]
The block record maps are jointly measurable: their graph and assignment coordinates are finite, and their outcome coordinates are affine functions of the baselines for each fixed partition and assignment. Since \(h=\kappa s\), \cref{obj:thm:block-testing-scale} gives
\[
\operatorname{TV}(P_{+1,h},P_{-1,h})\le\frac12.
\]
Therefore \cref{obj:lem:full-record-two-prior-risk} yields
\[
R_{n,1}(q)
\ge
\frac{\Delta^2}{2}
\left(1-\operatorname{TV}(P_{+1,h},P_{-1,h})\right)
\ge
\frac{\Delta^2}{4}
=
\frac14\left(\frac{m\kappa s}{n}\right)^2.
\]
% lean: CausalSmith.Experimentation.ThinnedgraphAdditiveRiskFrontier.block_testing_minimax_lower

Finally, the rounding bound gives
\[
\frac mn\ge\frac14,
\qquad
\frac{\kappa s}{4}
\le\frac{m\kappa s}{n}.
\]
Both sides of the latter inequality are positive, so
\[
\left(\frac{\kappa s}{4}\right)^2
\le
\left(\frac{m\kappa s}{n}\right)^2,
\qquad
\frac14\left(\frac{\kappa s}{4}\right)^2
=
\frac{s^2}{640000\pi^2}.
\]
Combining these estimates with \(F(n,1,q)\le s^2\) gives
\[
\begin{aligned}
cF(n,1,q)
&\le \frac{s^2}{2560000\pi^2}\\
&\le \frac{s^2}{640000\pi^2}\\
&=\frac14\left(\frac{\kappa s}{4}\right)^2\\
&\le\frac14\left(\frac{m\kappa s}{n}\right)^2\\
&\le R_{n,1}(q).
\end{aligned}
\]
Together with the first case, this proves the degree-one frontier bound.
% lean: CausalSmith.Experimentation.ThinnedgraphAdditiveRiskFrontier.frontier_minimax_lower

\item The model class of \cref{obj:def:model} is increasing in its degree bound:
\[
\mathcal M_{n,1}\subseteq\mathcal M_{n,d}
\qquad(d\ge1).
\]
Indeed, the coefficient-mass condition is the same, and each degree bound at most one is also at most \(d\). For every measurable randomized estimator \(T\),
\[
\sup_{\theta\in\mathcal M_{n,1}}\mathcal L(T;\theta,q)
\le
\sup_{\theta\in\mathcal M_{n,d}}\mathcal L(T;\theta,q).
\]
Taking infima over the common estimator class gives
\[
R_{n,1}(q)\le R_{n,d}(q).
\]
% lean: CausalSmith.Experimentation.ThinnedgraphAdditiveRiskFrontier.minimaxRisk_mono_degree

We also have
\[
\frac1{1+nq}\le F(n,1,q).
\]
For \(q>0\), \(nq>0\), and
\[
\frac1{1+nq}\le1,
\qquad
\frac1{1+nq}\le\frac1{nq}.
\]
These two inequalities imply the claim by the degree-one formula for \(F\). At \(q=0\), both sides equal one. Multiplication by \(c>0\), followed by the bounds just proved, therefore gives
\[
\frac{c}{1+nq}
\le cF(n,1,q)
\le R_{n,1}(q)
\le R_{n,d}(q).
\]
% lean: CausalSmith.Experimentation.ThinnedgraphAdditiveRiskFrontier.degree_one_frontier_reciprocal

\item Combining the two lower bounds and using \(c>0\), we obtain
\[
\begin{aligned}
R_{n,d}(q)
&\ge
\max\left\{
c\min\left\{1,\frac{(d+1)^2}{n}\right\},
\frac{c}{1+nq}
\right\}\\
&=
c\max\left\{
\min\left\{1,\frac{(d+1)^2}{n}\right\},
\frac1{1+nq}
\right\}.
\end{aligned}
\]
% lean: CausalSmith.Experimentation.ThinnedgraphAdditiveRiskFrontier.supported_boundaries

\item At \(q=0\), the retention lower bound and the upper bound from \cref{obj:thm:observable-upper} give
\[
c\le R_{n,d}(0)\le1.
\]
% lean: CausalSmith.Experimentation.ThinnedgraphAdditiveRiskFrontier.supported_boundaries

\item At \(q=1\), the audit contribution in \cref{obj:def:upper-envelope} vanishes:
\[
u(n,d,1)=\frac{5(d+1)^2}{n}.
\]
Thus
\[
R_{n,d}(1)\le
\min\left\{1,\frac{5(d+1)^2}{n}\right\}.
\]
If \((d+1)^2/n\le1\), then
\[
\min\left\{1,\frac{5(d+1)^2}{n}\right\}
\le\frac{5(d+1)^2}{n}
\le
21\min\left\{1,\frac{(d+1)^2}{n}\right\}.
\]
If \((d+1)^2/n>1\), then
\[
\min\left\{1,\frac{5(d+1)^2}{n}\right\}
\le1\le21
=
21\min\left\{1,\frac{(d+1)^2}{n}\right\}.
\]
Together with the degree lower bound, these estimates prove
\[
c\min\left\{1,\frac{(d+1)^2}{n}\right\}
\le R_{n,d}(1)
\le
C_0\min\left\{1,\frac{(d+1)^2}{n}\right\}.
\]
% lean: CausalSmith.Experimentation.ThinnedgraphAdditiveRiskFrontier.supported_boundaries

\item It remains to bound the degree-one upper envelope. Suppose first that \(q>0\). Since \(q\le1\),
\[
u(n,1,q)
=
\frac{20}{n}+\frac{4(1-q)}{nq}
=
\frac{4+16q}{nq}
\le\frac{20}{nq}.
\]
Consequently,
\[
\min\{1,u(n,1,q)\}
\le
\min\left\{1,\frac{20}{nq}\right\}.
\]
If \(nq\le20\), then
\[
\min\left\{1,\frac{20}{nq}\right\}
\le1\le\frac{21}{1+nq}.
\]
If \(nq>20\), then the positive denominators and
\[
20(1+nq)\le21nq
\]
give
\[
\min\left\{1,\frac{20}{nq}\right\}
\le\frac{20}{nq}
\le\frac{21}{1+nq}.
\]
At \(q=0\), the envelope is \(+\infty\), so
\[
\min\{1,u(n,1,0)\}=1\le21=\frac{21}{1+nq}.
\]
Hence, throughout \(q\in[0,1]\),
\[
\min\{1,u(n,1,q)\}\le\frac{21}{1+nq}.
\]
Applying the upper bound from \cref{obj:thm:observable-upper} and the retention lower bound already obtained proves
\[
\frac{c}{1+nq}
\le R_{n,1}(q)
\le\frac{C_0}{1+nq}.
\]
% lean: CausalSmith.Experimentation.ThinnedgraphAdditiveRiskFrontier.degree_one_envelope_upper
\end{enumerate}
\end{proof}
